% TOP_PROBLEM: {"problemId": "problem.hilberts-tenth-problem-over-the-rationals", "canonicalTitle": "Hilbert's Tenth Problem over the Rationals", "releaseRank": 21, "primaryDomain": "logic_foundations_set_theory", "primaryDomainLabel": "Logic, foundations & set theory", "displayStatus": "open", "releaseStatus": "open"}
% !TeX program = lualatex
% ATTEMPT_STATUS: unresolved
\documentclass[11pt]{article}
\usepackage[margin=25mm]{geometry}
\usepackage{fontspec}
\setmainfont{Latin Modern Roman}
\usepackage{amsmath,amssymb,mathtools}
\usepackage{unicode-math}
\setmathfont{Latin Modern Math}
\usepackage{xurl,hyperref}
\hypersetup{colorlinks=true,urlcolor=blue}
\setcounter{secnumdepth}{0}
\setlength{\parindent}{0pt}
\setlength{\parskip}{5pt}
\setlength{\emergencystretch}{3em}
\begin{document}
\section*{\texorpdfstring{Hilbert's Tenth Problem over the Rationals}{Hilbert's Tenth Problem over the Rationals}}\label{21}
\subsection{Definitions and mathematical statement}
An input is a finite description of $f\in{\mathbb{Z}}[x_1,\ldots,x_n]$, with no fixed bound on $n$, degree, or coefficient size. Define the decision set
\[
 \mathcal D_{{\mathbb{Q}}}=\{\langle f\rangle:\exists(q_1,\ldots,q_n)\in{\mathbb{Q}}^n,
                                  f(q_1,\ldots,q_n)=0\}.
\]
Determine whether $\mathcal D_{{\mathbb{Q}}}$ is decidable: one algorithm must halt on every input and return its correct membership bit. Enumeration of rational solutions only semidecides yes-instances and is not a solution. The selected assertion concerns computability, with no polynomial-time demand.
\par\smallskip\textit{Formulation and concept references:} \href{https://arxiv.org/pdf/2309.14987}{[S1]}.

\subsection{Short English statement}
Can an algorithm always decide whether a polynomial equation with integer coefficients has a rational solution?
\subsection{Research attempt}
This unresolved attempt by GPT-6 Astra Ultra was prepared on 27 September
2026, using explicit reductions, source checks, and adversarial verification.

\paragraph{Scope, source check, and attack map.}
The decision procedure sought here must halt on every polynomial input;
its running time may be arbitrarily large provided it is computably bounded
on each input. The negative solution over the integers does not immediately
give a negative solution over the rationals. Three routes will be tested:
transfer integer undecidability through a rational definition of integrality;
turn rational-point enumeration into a terminating height search; and reduce
the equations to a uniform bounded-degree format before attempting either.

Two inherited references require particular care. Koymans--Pagano \href{https://arxiv.org/html/2602.04468v1}{[S2]},
Theorem 1.1 and Question 6.2, treat infinite finitely generated rings and
explicitly leave $\mathbb Q$ outside that result. Garcia-Fritz--Pasten--Vidaux
\href{https://arxiv.org/abs/2311.01958}{[S3]} prove undecidability with extra height-comparison predicates, which are
not part of the present input language. Koenigsmann \href{https://arxiv.org/abs/1011.3424}{[E1]} gives a universal
definition of $\mathbb Z$ in $\mathbb Q$; this is not an existential definition.
None of these facts is used with its qualifiers removed. The constructions
below are elementary reductions, not claims of new undecidability theorems.

\paragraph{First test: clearing denominators does not transfer the answer.}
For $f\in\mathbb Z[x_1,\ldots,x_n]$ of total degree $d$, define
\[
 F(a_1,\ldots,a_n,b)=b^df(a_1/b,\ldots,a_n/b).
\]
This is an integer polynomial, and $f$ has a rational zero exactly when
$F$ has an integer zero with $b\ne0$. The forward direction uses a common
positive denominator; the reverse direction divides by a nonzero $b$.
The condition $b\ne0$ is essential. For instance, the homogenization of
$x^2+y^2+1$ has the zero vector as an integer zero, although the original
polynomial is positive on every rational pair. Its zero at the origin of
the homogeneous coordinate space is not a rational affine point.

This implication reduces the rational question to a special collection
of integer existence questions. Undecidability of the larger collection
does not imply undecidability of the special one. In the other direction,
an integer polynomial may acquire rational solutions absent over the
integers: $2x-1=0$ is the simplest example. Thus simply using the same
polynomial does not reduce integer solvability to rational solvability.
The missing operation in that proposed reduction is restriction of the
quantified variables to integers.

The finite-generation distinction can also be checked without a theorem.
For rationals $r_1,\ldots,r_s$, let $D$ be a common multiple of their
denominators. Then
\[
 \mathbb Z[r_1,\ldots,r_s]\subseteq\mathbb Z[1/D].
\]
Every denominator on the right has all prime factors dividing $D$.
If $p\nmid D$ is prime, $1/p$ is missing. Hence $\mathbb Q$ is not
finitely generated as a ring. Taking a union over all finite prime sets
does not turn a theorem about any one of these rings into a decision
procedure, or an undecidability reduction, over their union.

\paragraph{Lemma 21.1: the exact definability bridge.}
Suppose there is a fixed integer polynomial $P(t,u_1,\ldots,u_m)$ such that
\begin{equation}\label{cont21:integrality}
 t\in\mathbb Z\quad\Longleftrightarrow\quad
 \exists u\in\mathbb Q^m:\ P(t,u)=0
 \qquad(t\in\mathbb Q).
\end{equation}
Then Hilbert's tenth problem over $\mathbb Q$ is undecidable.

\emph{Proof.} Given any $f\in\mathbb Z[x_1,\ldots,x_n]$, introduce
separate witness tuples $u^{(1)},\ldots,u^{(n)}$ and form
\[
 Q_f(x,u)=f(x)^2+\sum_{i=1}^n P(x_i,u^{(i)})^2.
\]
Over $\mathbb Q\subset\mathbb R$, a sum of squares vanishes if and only
if every summand vanishes. Thus $Q_f$ has a rational zero precisely when
$f$ has an integer zero. The map $f\mapsto Q_f$ is effective because
$P$ is fixed. A rational-solvability decider would therefore decide the
integer problem, contrary to the Matiyasevich--Robinson--Davis--Putnam
theorem recalled in \href{https://arxiv.org/pdf/2309.14987}{[S1]}.\hfill$\square$

The hypothesis is not established here. Replacing it by
$t\in\mathbb Z\Longleftrightarrow\forall u\,\Phi(t,u)$ changes the
resulting formula to one with universal quantifiers after the existential
$x$ variables. The present decision oracle answers only existential
polynomial questions. Negating the universal formula provides witnesses
for nonintegrality, not witnesses for integrality; negation cannot simply
be moved through $\exists x$ while preserving the problem. Consequently
Koenigsmann's theorem does not supply the displayed $P$ by a syntactic
change of quantifier. A Diophantine model of integer arithmetic would be
another sufficient route, but none is constructed in this attempt.

\paragraph{Lemma 21.2: degree four already contains the full difficulty.}
Every rational polynomial-solvability question can be effectively converted
to one equation of total degree at most four, allowing additional variables.

\emph{Proof.} Evaluate the input polynomial by a finite arithmetic circuit
with addition, multiplication and integer constants. Repeated squaring
handles exponents encoded in binary. Attach a rational variable $y_j$ to
each gate. An addition gate is enforced by
$g_j=y_j-y_a-y_b=0$, a multiplication gate by
$g_j=y_j-y_ay_b=0$, and a constant gate by $g_j=y_j-c=0$.
Include the output gate equation $y_{\mathrm{out}}=0$. Form
\begin{equation}\label{cont21:quartic}
 Q(x,y)=y_{\mathrm{out}}^2+\sum_j g_j(x,y)^2.
\end{equation}
Each $g_j$ has degree at most two, so $Q$ has degree at most four.
If $f(x)=0$, the actual gate values give a rational zero of $Q$.
Conversely, a rational zero forces each gate equation to hold; induction
in circuit order shows that $y_{\mathrm{out}}=f(x)=0$.
The construction terminates and writes an explicit integer polynomial.
\hfill$\square$

Constants and identically zero polynomials can be handled directly; free
unused input variables cause no problem. The ordering of $\mathbb Q$ is
used in the sum-of-squares step, so this argument is not asserted over
an arbitrary field. The number of variables grows with the circuit.
Thus a solution for quartics with unbounded dimension would solve the
whole target; a result for one fixed number of variables would not do so
without a further reduction.

\paragraph{Attempt B: replace enumeration by a certified height bound.}
For a reduced fraction $a/b$, $b>0$, define
$H(a/b)=\max\{|a|,b\}$ and
$H(q_1,\ldots,q_n)=\max_iH(q_i)$, with height one for an empty tuple.
Each bounded-height box is finite and effectively enumerable. Exact
rational arithmetic decides whether any given tuple is a zero.
Increasing the bound therefore semidecides positive instances.

\textbf{Proposition 21.3.} Fix any effective finite binary encoding of the
polynomial inputs, of length $L(f)$. The following are equivalent:
(i) the rational-solvability set is decidable;
(ii) there is a total computable $B:\mathbb N\to\mathbb N$ such that
every solvable $f$ has a zero of height at most $B(L(f))$.

\emph{Proof.} Given (ii), compute $B(L(f))$, enumerate the finite box,
and test every tuple. Return no if none is a zero. This halts and is
correct by the bound. For the reverse implication, enumerate all valid
input strings of length at most $L$ and use the hypothetical decider on
each. There are finitely many. For every one declared solvable, run the
rational enumeration until its first solution is found. Each such search
halts because its positive decision was correct. Return the maximum of
their heights and one. This is a terminating algorithm for $B(L)$ and
has the required property.\hfill$\square$

This is a proved reduction, not an estimate for $B$. It isolates the
effective ingredient missing from brute force. Mere finiteness of the
set of inputs at each length gives a mathematical maximum of least
solution heights, but it does not establish that the maximum is
computable. Conversely, no polynomial, exponential, or otherwise modest
growth is required of $B$ for decidability. Showing that one proposed
small bound fails cannot disprove existence of every computable bound.

\paragraph{An exact quartic stress test for small-height guesses.}
For $m\ge1$, consider
\begin{equation}\label{cont21:chain}
 Q_m(x_1,\ldots,x_m)=(x_1-2)^2
                  +\sum_{i=1}^{m-1}(x_{i+1}-x_i^2)^2.
\end{equation}
The nonnegative summands force the unique rational zero
\[
 x_i=2^{2^{i-1}}\quad(1\le i\le m),\qquad
 H(x)=2^{2^{m-1}}.
\]
For $m\ge2$ the degree is four, the coefficients after expansion are
bounded independently of $m$, and there are $O(m)$ monomials. In a sparse
encoding with binary variable indices the input length is
$O(m\log(m+2))$. The largest numerator needs $2^{m-1}+1$ binary digits.
It therefore defeats every bound polynomial in input length on the bit
length of a required witness. For $m=1$ the equation simply has the
unique zero $2$, agreeing with the formula.

This is an obstruction to a small-witness shortcut, not an undecidability
proof: the displayed double-exponential height is entirely computable.
It also rules out a height bound depending only on degree and the largest
coefficient while ignoring the number of variables. The quantifier over
dimension in the original question is doing real work.

\paragraph{Why finite local checks are not a general negative certificate.}
Even checking integrality of one rational against a fixed list of primes
does not enforce integrality at every prime. For any finite set $S$, choose
a prime $r\notin S$. Then $1/r$ is integral in $\mathbb Q_p$ for each
$p\in S$, but is not an integer. This elementary test invalidates a
proposal that simply substitutes finitely many local integrality tests
for \eqref{cont21:integrality}. It does not exclude more sophisticated
finite existential encodings using additional variables. Likewise,
local-global theorems for special polynomial families require their
hypotheses; no general local-global principle is inferred from passing
any finite list of congruence tests.

\paragraph{Verification, first gaps, and outcome.}
The companion exact-arithmetic checks evaluate the squaring chains,
confirm their heights, and test the gate-equation transformation on a
small rational grid. The proofs, rather than those finite samples,
establish the reductions for every polynomial. The original choice of
rational witnesses, all dimensions, and unrestricted computable running
time is retained throughout.

\textbf{UNRESOLVED.} The positive route stops at a total computable height
bound, and the negative route stops at an existential arithmetic model
or a definition such as \eqref{cont21:integrality}. The quartic reduction
and the explicit height family identify constraints any continuation
must respect. Next steps are to construct the missing existential model,
to obtain certified height bounds on a larger geometric class, or to
prove that a proposed local-global algorithm supplies both positive and
negative termination certificates. No universal algorithm or proof of
its impossibility is obtained here.

\subsection{Sources}
\begin{itemize}
\item[S1]Sylvy Anscombe, Valentijn Karemaker, Zeynep Kisakürek, Vlerë Mehmeti, Margherita Pagano and Laura Paladino, A survey of local-global methods for Hilbert's Tenth Problem, abstract and Section1.
\url{https://arxiv.org/pdf/2309.14987}.
\textit{Formulation source.}
\item[S2]Hilbert's tenth problem for finitely generated rings.
\url{https://arxiv.org/html/2602.04468v1}.
\textit{Further reference; not a proof endorsement.}
\item[S3]Effectivity for existence of rational points is undecidable.
\url{https://arxiv.org/abs/2311.01958}.
\textit{Further reference; not a proof endorsement.}
\item[E1]Jochen Koenigsmann, \emph{Defining $\mathbb Z$ in $\mathbb Q$},
Annals of Mathematics 183 (2016), 73--93; arXiv:1011.3424.
\url{https://arxiv.org/abs/1011.3424}.
\textit{Universal definition, not the existential definition assumed in
Lemma 21.1; consulted 27 September 2026.}
\end{itemize}
\end{document}
